\documentclass[11pt,a4paper]{article}
\usepackage[T1]{fontenc}
\usepackage[utf8]{inputenc}
\usepackage{lmodern}
\usepackage[margin=20mm]{geometry}
\usepackage{amsmath,amssymb}
\usepackage{enumitem}
\usepackage{xcolor}
\usepackage{hyperref}
\usepackage{fancyhdr}
\usepackage{mdframed}

\definecolor{linkblue}{HTML}{245579}
\definecolor{replyblue}{HTML}{356FA3}
\definecolor{replyfill}{HTML}{F1F6FB}
\newmdenv[
  linewidth=0.6pt,
  linecolor=replyblue,
  backgroundcolor=replyfill,
  innerleftmargin=8pt,
  innerrightmargin=8pt,
  innertopmargin=6pt,
  innerbottommargin=6pt,
  skipabove=0.6em,
  skipbelow=0.8em,
  frametitlefont=\bfseries\small,
  font=\small,
  repeatframetitle=true,
  nobreak=false
]{responsebox}
\hypersetup{
  colorlinks=true,
  linkcolor=linkblue,
  urlcolor=linkblue,
  pdftitle={Beyond Imitation -- Response to Reviewers},
  pdfauthor={}
}
\setlength{\parindent}{0pt}
\setlength{\parskip}{0.2em}
\setlist{leftmargin=*,topsep=0.2em,itemsep=0.2em,parsep=0pt}
\setlength{\emergencystretch}{2em}
\setlength{\headheight}{14pt}
\pagestyle{fancy}
\fancyhf{}
\fancyhead[L]{\small Beyond Imitation}
\fancyhead[R]{\small Responses to Reviewers}
\fancyfoot[C]{\thepage}
\renewcommand{\headrulewidth}{0.3pt}
\raggedbottom

\title{\Large Beyond Imitation: Filtering On-Policy Distillation\\
by Reasoning Progress\\[0.65em]
\LARGE Response to Reviewers}
\author{}
\date{}

\begin{document}
\maketitle
\thispagestyle{fancy}

\textbf{AAAI 2027 -- Submission 11993}

The five official human reviews are reproduced in their order in the supplied
OpenReview PDF. Reviewer numbers are local labels for this working document;
the original anonymous reviewer IDs are retained below.
The separate AI review is included in Appendix~\ref{reviewer:ai}.
Blue boxes contain the author's existing draft responses or space for further
drafting. Every weakness and rebuttal question has a corresponding response box.
Each review ends with an overall response box. Unfilled boxes are marked
\textit{To be completed.}

\begin{center}
\begin{tabular}{lccc}
\hline
Reviewer & OpenReview ID & Rating & Confidence \\
\hline
Reviewer 1 & \texttt{gURK} & 5 & 4 \\
Reviewer 2 & \texttt{4R1c} & 5 & 4 \\
Reviewer 3 & \texttt{XU6y} & 5 & 4 \\
Reviewer 4 & \texttt{hTMj} & 2 & 4 \\
Reviewer 5 & \texttt{YgpD} & 6 & 4 \\
\hline
\end{tabular}
\end{center}

\tableofcontents

% Each reviewer has a separate file for collaborative editing.
\clearpage

% Inlined snapshot: reviewers/reviewer1
% Reviewer 1; OpenReview ID: gURK
% Original review text, followed by the author's supplied draft responses.
\section{Reviewer 1 (\texttt{gURK})}
\label{reviewer:1}
\noindent\textbf{Rating:} 5: Marginally below acceptance threshold\par
\noindent\textbf{Confidence:} 4: The reviewer is confident but not absolutely certain that the evaluation is correct

\subsection*{Paper Summary}
This paper starts from the argument that in OPD the student model's per token level agreement with teacher model is not perfect way for improving its reasoning process, as a student's generation can be moving towards the correct solution while having a high KL due to a different distribution from the teacher's. This paper then presents R\textsuperscript{2}-OPD, which uses a teacher-independent progress signal only as a reliability filter on the distillation loss. By splitting the student model's thinking trajectories into multiple segments by a preset markers, and process the reward with segments that contributes to the final ground truth answer, this method improves reasoning-oriented OPD, especially on harder benchmarks.

\subsection*{Strengths}
\begin{itemize}
\item This paper is well motivated and generally well written with a clear separation from distillation compatibility from reasoning progress, which targets the fact that OPD treats the former as a proxy for the latter and thus impacts the training efficiency.
\item The method is straightforward on the design including segmentation methods and easy to adopt for other methods. Specifically on the segmentation, the progress signal acts as a gate when applies the rules, and falls back to OPD when not available, the conflicts inside response are also addressed.
\item The experiments are well designed and clearly demonstrated, with Fig 1 showing student rollouts on DAPO that contain the answer but go unscored, which is a concrete argument that outcome-only filtering discards usable trajectories.On Fig 3-5 the paper also shows the configurations and ablations over segments merging strategy etc.
\end{itemize}

\subsection*{Weaknesses}
\begin{itemize}
\item This paper seems a combination of several existing works and some new components are not shown to be necessary through ablations like process reward or the process of filtering OPD supervision with the teacher-independent signal. It seems the genuine new components are mostly segmentation methods, conflict resolving rule and top-q masking, which are mostly heuristic and not compared with similar or simpler methods.
\par\Needspace{6\baselineskip}
\begin{responsebox}[frametitle={Response to W1}]
% Edit the author response below. Reviewer text above is unchanged.
The key contribution of R\textsuperscript{2}-OPD is a novel progress-aware criterion that identifies unreliable teacher supervision through ranking conflicts between independently estimated reasoning progress and distillation loss, rather than relying on token-level heuristics or outcome correctness alone. Sign-consistent merging and selective masking operationalize this criterion to suppress conflicting supervision while preserving useful guidance. We will clarify this distinction and acknowledge the need for further component-level ablations.
\end{responsebox}

\item The overall performance imprvements seem limited and specific to one benchmark family. In table 2 the gain is +0.43 avg@4, and over AIME25 avg@4 decreases.
\par\Needspace{6\baselineskip}
\begin{responsebox}[frametitle={Response to W2}]
% Edit the author response below. Reviewer text above is unchanged.
GPQA-Diamond: avg@4 / pass@4 (\%)

\begin{center}
\begin{tabular}{lccc}
\hline
Model & Student & OPD & R\textsuperscript{2}-OPD \\
\hline
Qwen2.5 & 33.46/67.17 & 35.23/67.68 & 36.24/69.19 \\
Qwen3 & 36.62/59.09 & 39.52/66.16 & 41.04/65.66 \\
\hline
\end{tabular}
\end{center}
\end{responsebox}

\item The evaluation design and results can be better explained. First of as the paper also suggests, only a small fraction of DAPO rollouts are scored, training at 7k tokens and evaluating at 8k might be truncating R1-Distill-1.5B responses as well, so shortening the responses will naturally improve the performance. Also it seems most baselines E-OPD, TIP-OPD, IW-OPD and Uni-OPD all score below the untrained student on AIME24, it might indicates the specific configuration favors the method in the paper. Moreover, the teacher model JustRL-DeepSeek-1.5B, is based on DeepSeek-R1-Distill-Qwen-1.5B, which comes with a low-divergence behavior and does not really reflect the motivation of this paper.
\par\Needspace{6\baselineskip}
\begin{responsebox}[frametitle={Response to W3}]
% Edit the author response below. Reviewer text above is unchanged.
We hypothesize that the performance gap partly reflects differences in reasoning regimes. E-OPD and TIP-OPD preserve or prioritize uncertain tokens, which may encourage additional exploration without improving answer quality in thinking-enabled models. IW-OPD downweights later tokens based on accumulated teacher--student discrepancy, potentially suppressing useful verification and correction in long reasoning traces. These effects may be amplified by our response-length limit. Our teacher selection follows arXiv:2604.13016, which finds that effective OPD requires teachers to provide new reasoning capabilities acquired through additional RL post-training, rather than simply having more parameters.
\end{responsebox}

\item The theory part is thin, Lemma 1 holds for fixed runs only which is natually too idealized. With prop.2, neither the 2 really analyzes the main contribution of the segmentation or masking strategy.
\par\Needspace{6\baselineskip}
\begin{responsebox}[frametitle={Response to W4}]
% Edit the author response below. Reviewer text above is unchanged.
Lemma 1 and Proposition 2 jointly support the stability of process-reward and KL estimates, enabling more reliable comparisons for conflict detection and masking. Our masking strategy is inspired by existing OPD methods: TIP-OPD and IW-OPD explore selective token-level weighting, while Uni-OPD filters supervision at the trajectory level. We extend this idea to the reasoning-span level, using progress--distillation conflicts to guide selective masking.
\end{responsebox}

\end{itemize}

\subsection*{Justification of Recommendation}
Based on the 3 weakness over limited performance improvements and evaluation design, I'm leaning towards a boarderline rejection.

\subsection*{Specific Points of Feedback for Rebuttal}
In addition to the weakness,

\begin{itemize}
\item The evaluation results could benefit from avg@32 on AIME with different seeds, more benchmarks.
\par\Needspace{6\baselineskip}
\begin{responsebox}[frametitle={Response to Q1},nobreak=true]
% Edit the author response below. Reviewer text above is unchanged.
\textit{[To be completed.]}\par\vspace{1em}
\end{responsebox}

\item Clear demonstration of segmentation and masking, e.g. the fraction of masked tokens that occur after the ground-truth answer first appears
\par\Needspace{6\baselineskip}
\begin{responsebox}[frametitle={Response to Q2},nobreak=true]
% Edit the author response below. Reviewer text above is unchanged.
\textit{[To be completed.]}\par\vspace{1em}
\end{responsebox}

\item Response length and truncation rate at 8k for student models over baselines
\par\Needspace{6\baselineskip}
\begin{responsebox}[frametitle={Response to Q3},nobreak=true]
% Edit the author response below. Reviewer text above is unchanged.
\textit{[To be completed.]}\par\vspace{1em}
\end{responsebox}

\item Justification of hyperparameters and implementation method of the baselines
\par\Needspace{6\baselineskip}
\begin{responsebox}[frametitle={Response to Q4},nobreak=true]
% Edit the author response below. Reviewer text above is unchanged.
\textit{[To be completed.]}\par\vspace{1em}
\end{responsebox}

\end{itemize}

\subsection*{Overall Response}
\par\Needspace{6\baselineskip}
\begin{responsebox}[frametitle={Overall Response to Reviewer 1 (gURK)},nobreak=true]
% Edit the author response below. Reviewer text above is unchanged.
\textit{[To be completed.]}\par\vspace{1em}
\end{responsebox}



\clearpage

% Inlined snapshot: reviewers/reviewer2
% Reviewer 2; OpenReview ID: 4R1c
% Original review text, followed by the author's supplied draft responses.
\section{Reviewer 2 (\texttt{4R1c})}
\label{reviewer:2}
\noindent\textbf{Rating:} 5: Marginally below acceptance threshold\par
\noindent\textbf{Confidence:} 4: The reviewer is confident but not absolutely certain that the evaluation is correct

\subsection*{Paper Summary}
This paper argues that OPD's use of teacher similarity as a proxy for reasoning progress is unreliable, as teacher-derived rewards may conflict with actual reasoning progress. The authors propose R\textsuperscript{2}-OPD, which estimates process rewards from on-policy rollouts, merges sign-consistent spans, and compares process-reward rankings with segment-level KL rankings. The top-q segments with the largest ranking conflicts are then masked.

\subsection*{Strengths}
I agree that this is a meaningful problem: dense token-level supervision in OPD does not necessarily reflect reasoning progress. Figure 2 presents a concrete and convincing failure mode, where regressive segments actually have lower KL divergence. Intuitively, using sampling-based process rewards as a teacher-independent reference to identify and mask unreliable teacher signals seems reasonable.

\subsection*{Weaknesses}
\begin{enumerate}
\item Sec. 4.1 reports $(N_{\mathrm{eval}}=8)$ and a maximum of 300 tokens per rollout, but does not quantify the additional generation overhead per training step relative to standard OPD. It would be helpful to report the total rollout or generated-token budget and its ratio to standard OPD.
\par\Needspace{6\baselineskip}
\begin{responsebox}[frametitle={Response to W1}]
% Edit the author response below. Reviewer text above is unchanged.
We sincerely thank the reviewer. Due to time constraints, we analyzed six sampled batches. PR was evaluated for only 103/192 responses (53.6\%), as responses without successful segmentation or explicit ground-truth answer matches bypass PR estimation and remain unfiltered. Eligible responses averaged 214.8 short continuations (eight per non-final prefix, each capped at 300 new tokens), yielding 115.7 additional rollouts per original response overall, including padding. PR estimation and consistency checking incurred 118.4\% overhead relative to the inferred step-time baseline. We will report this breakdown and acknowledge the lack of compute-matched comparisons.
\end{responsebox}

\item My main concern is scalability and generalizability. The experiments only use DeepSeek-R1-Distill-Qwen-1.5B, and the training dataset \& evaluation also appear limited. This makes it difficult to assess whether the proposed method scales to larger or more recent models.
\par\Needspace{6\baselineskip}
\begin{responsebox}[frametitle={Response to W2},nobreak=true]
% Edit the author response below. Reviewer text above is unchanged.
\textit{[To be completed.]}\par\vspace{1em}
\end{responsebox}

\item On AIME 25, performance drops sharply from 25.83 to 12.15 when (q=50). Such a large degradation suggests that the optimal (q) may strongly depend on the noise level of the process-reward estimates, raising concerns about hyperparameter sensitivity.
\par\Needspace{6\baselineskip}
\begin{responsebox}[frametitle={Response to W3}]
% Edit the author response below. Reviewer text above is unchanged.
We thank the reviewer for this observation. We acknowledge the sensitivity to q, which controls the fraction of masked segments rather than tokens. Aggressive masking may remove useful supervision or amplify ranking errors, but the degradation cannot be attributed solely to PR noise. We will clarify this trade-off and explore token-level masking budgets or soft weighting to improve robustness.
\end{responsebox}

\item Unmasking OPD (arXiv:2605.10889) shows through token-level gradient alignment analysis that distillation guidance on correct rollouts "tends to become noisy." R\textsuperscript{2}-OPD uses rank conflicts ($\kappa_j$) to identify unreliable teacher signals, but the findings of Unmasking OPD suggest that such unreliability may be concentrated in later segments of correct trajectories rather than uniformly distributed. How do the authors view this observation in relation to their rank-conflict criterion?
\par\Needspace{6\baselineskip}
\begin{responsebox}[frametitle={Response to W4}]
% Edit the author response below. Reviewer text above is unchanged.
We thank the reviewer for this insightful observation. To investigate this, we conducted exploratory offline replay on 8 saved batches (256 trajectories). Among the 39 correctly answered trajectories with masked spans, 59.8\% of masked tokens occurred before the final third of the response, while the overlap with a late-only mask under an equal token budget was only 19.6\%. These results suggest that R\textsuperscript{2}-OPD does not simply truncate supervision in later positions. We agree that position can influence signal reliability; however, our rank-conflict criterion identifies local disagreements between reasoning progress and distillation loss, rather than assuming that unreliability is uniformly distributed or intrinsically tied to position. We will clarify this distinction while noting that masking is an operational proxy, not direct ground-truth evidence of signal unreliability.
\end{responsebox}

\item Prune-OPD and Teachability-Aware OPD have also, to some extent, explored reallocating distillation budgets based on signal reliability. Although their implementations differ from R\textsuperscript{2}-OPD, discussing these related approaches would help clarify the novelty and positioning of the proposed method.
\par\Needspace{6\baselineskip}
\begin{responsebox}[frametitle={Response to W5}]
% Edit the author response below. Reviewer text above is unchanged.
Thanks. Prune-OPD detects prefix drift and reallocates the training budget through reward down-weighting and rollout truncation, while Teachability-Aware OPD selects token positions according to local teacher--student compatibility and teachability. In contrast, R\textsuperscript{2}-OPD uses independently estimated reasoning progress to identify rank conflicts at the reasoning-span level, selectively masking conflicting supervision while retaining the original OPD objective. These methods share the motivation of selective supervision but employ different reliability criteria and intervention mechanisms.

We will strengthen the related-work discussion to better articulate these distinctions and clarify that R\textsuperscript{2}-OPD complements, rather than replaces, existing reliability-aware approaches.
\end{responsebox}

\end{enumerate}

\subsection*{Justification of Recommendation}
meaningful problem yet constrained by the outdated model \& dataset selection

\subsection*{Specific Points of Feedback for Rebuttal}
See Weaknesses 1, 2, 3. Happy to raise the score given more persuative experiments.

\par\Needspace{6\baselineskip}
\begin{responsebox}[frametitle={Response to Q1},nobreak=true]
% Edit the author response below. Reviewer text above is unchanged.
\textit{[To be completed.]}\par\vspace{1em}
\end{responsebox}

\subsection*{Overall Response}
\par\Needspace{6\baselineskip}
\begin{responsebox}[frametitle={Overall Response to Reviewer 2 (4R1c)},nobreak=true]
% Edit the author response below. Reviewer text above is unchanged.
\textit{[To be completed.]}\par\vspace{1em}
\end{responsebox}



\clearpage

% Inlined snapshot: reviewers/reviewer3
% Reviewer 3; OpenReview ID: XU6y
% Original review text, followed by the author's supplied draft responses.
\section{Reviewer 3 (\texttt{XU6y})}
\label{reviewer:3}
\noindent\textbf{Rating:} 5: Marginally below acceptance threshold\par
\noindent\textbf{Confidence:} 4: The reviewer is confident but not absolutely certain that the evaluation is correct

\subsection*{Paper Summary}
The paper targets the improvement of the on-policy distillation (OPD) method. Specifically, the paper proposes R\textsuperscript{2}-OPD, a reasoning-progress-aware filtering method for on-policy distillation. R\textsuperscript{2}-OPD estimates the progress made by intermediate reasoning segments through Monte Carlo continuation rollouts and verifier-based solve probabilities, and then compares the within-trajectory ranking induced by estimated reasoning progress with that induced by the teacher-student divergence. When the two rankings conflict, the method masks the corresponding distillation loss, aiming to prevent teacher imitation from penalizing productive but teacher-divergent reasoning steps. Experiments on mathematrical reasoning demonstrate the effectiveness of the proposed method.

\subsection*{Strengths}
\begin{itemize}
\item The paper identifies a meaningful limitation of standard OPD, in which teacher supervision may not always align with genuine reasoning progress.
\item The idea is well-motivated and easy to understand.
\item The experiments include several relevant OPD baselines, including E-OPD, TIP-OPD, IW-OPD, and Uni-OPD.
\end{itemize}

\subsection*{Weaknesses}
\begin{itemize}
\item It seems that the proposed method introduces substantial computational overhead. Each eligible intermediate segment requires eight additional continuation rollouts, while responses can be thousands of tokens long. The paper does not report its training cost or provide compute-matched comparisons with other OPD methods.
\par\Needspace{6\baselineskip}
\begin{responsebox}[frametitle={Response to W1}]
% Edit the author response below. Reviewer text above is unchanged.
We sincerely thank the reviewer. Due to time constraints, we analyzed six sampled batches. PR was evaluated for only 103/192 responses, as responses without successful segmentation or explicit ground-truth answer matches bypass PR estimation and remain unfiltered. Eligible responses averaged 214.8 short continuations (eight per non-final prefix, each capped at 300 new tokens), yielding 115.7 additional rollouts per original response overall, including padding. PR estimation and consistency checking incurred 118.4\% overhead relative to the inferred step-time baseline. We will report this breakdown and acknowledge the lack of compute-matched comparisons.
\end{responsebox}

\item The proposed method requires annotated ground-truth answers and task-specific verifiers to estimate reasoning progress. It is therefore restricted to verifiable domains and may partly benefit from correctness information unavailable to standard OPD, which does not require ground-truth answers.
\par\Needspace{6\baselineskip}
\begin{responsebox}[frametitle={Response to W2}]
% Edit the author response below. Reviewer text above is unchanged.
We thank the reviewer for highlighting this limitation. While R\textsuperscript{2}-OPD requires ground-truth answers and task-specific verifiers, its contribution goes beyond incorporating correctness information. Unlike methods that directly use outcome correctness to filter or reweight trajectories, R\textsuperscript{2}-OPD estimates the marginal progress of intermediate reasoning spans and compares its relative ranking with that of distillation losses. This enables fine-grained identification of locally conflicting teacher signals, rather than merely distinguishing correct from incorrect responses. We acknowledge the additional supervision requirement and will clarify the applicability limitations and explore verifier-free progress estimation in future work.
\end{responsebox}

\item The paper provides limited analysis of training dynamics. It would be useful to show how reasoning-progress estimates, teacher--progress ranking conflicts, masking rates, and downstream performance evolve throughout training, to better understand the proposed idea and method.
\par\Needspace{6\baselineskip}
\begin{responsebox}[frametitle={Response to W3}]
% Edit the author response below. Reviewer text above is unchanged.
We thank the reviewer for this helpful suggestion. Figure 5 provides an initial view of training dynamics, where improved PR--KL ranking agreement with sign-consistent merging illustrates the role of reasoning-progress estimates in identifying teacher-signal conflicts. We agree that further analysis would be valuable. We note that the masking ratio (q=30\%) remains fixed throughout training.
\end{responsebox}

\end{itemize}

\subsection*{Justification of Recommendation}
I assign a score of 5, as the paper presents a well-motivated idea and shows performance gains over standard OPD. However, there are core concerns regarding its reliance on ground-truth answers and verifiers, potentially substantial but unreported computational overhead. Addressing these issues with clearer efficiency analysis and broader evaluations would strengthen the paper.

\subsection*{Specific Points of Feedback for Rebuttal}
See above weaknesses.

\par\Needspace{6\baselineskip}
\begin{responsebox}[frametitle={Response to Q1},nobreak=true]
% Edit the author response below. Reviewer text above is unchanged.
\textit{[To be completed.]}\par\vspace{1em}
\end{responsebox}

\subsection*{Overall Response}
\par\Needspace{6\baselineskip}
\begin{responsebox}[frametitle={Overall Response to Reviewer 3 (XU6y)},nobreak=true]
% Edit the author response below. Reviewer text above is unchanged.
\textit{[To be completed.]}\par\vspace{1em}
\end{responsebox}



\clearpage

% Inlined snapshot: reviewers/reviewer4
% Reviewer 4; OpenReview ID: hTMj
% Original review text, followed by the author's supplied draft responses.
\section{Reviewer 4 (\texttt{hTMj})}
\label{reviewer:4}
\noindent\textbf{Rating:} 2: Strong rejection\par
\noindent\textbf{Confidence:} 4: The reviewer is confident but not absolutely certain that the evaluation is correct

\subsection*{Paper Summary}
This paper proposes R2-OPD, a selective on-policy distillation method motivated by the potential mismatch between teacher agreement and reasoning progress. The method samples continuations from reasoning prefixes to estimate changes in the probability of obtaining a correct answer, merges adjacent spans with sign-consistent progress, and masks supervision where progress and average distillation-loss rankings conflict. Overall, the paper proposes a highly incremental design, and the scope of the experiments, the analysis, and the results are all quite limited.

\subsection*{Strengths}
\begin{enumerate}
\item Teacher agreement can differ from reasoning progress, making selective distillation a useful direction. Applying existing progress-based supervision [2] to span-level filtering is a concrete proposal.
\item Table 1 reports improvements over OPD in avg@4 and pass@4. The masking and merging ablations demonstrate sensitivity to these design choices, without establishing the claimed mechanism.
\end{enumerate}

\subsection*{Weaknesses}
\begin{enumerate}
\item The method combines existing ideas from progress-based supervision and reward-guided distillation [1,2]. Its main distinction is using estimated progress to select individual reasoning spans, which appears to be a relatively narrow methodological extension.
\par\Needspace{6\baselineskip}
\begin{responsebox}[frametitle={Response to W1},nobreak=true]
% Edit the author response below. Reviewer text above is unchanged.
\textit{[To be completed.]}\par\vspace{1em}
\end{responsebox}

\item Both evaluated students are small Qwen-derived models. This makes it difficult to tell whether the gains extend beyond these particular models. Results on another model family, such as Llama, and larger students would be needed to support broader claims. All evaluations use a single response-length limit. The paper does not show whether the gains remain when longer reasoning is allowed. RG-OPD is closely related to the proposed method but is not included in the experiments [1]. Comparisons with the closest existing methods and simpler reward-based filters are needed.
\par\Needspace{6\baselineskip}
\begin{responsebox}[frametitle={Response to W2},nobreak=true]
% Edit the author response below. Reviewer text above is unchanged.
\textit{[To be completed.]}\par\vspace{1em}
\end{responsebox}

\item The filter relies on progress estimates from continuations limited to 300 tokens. It is unclear whether these short, answer-eliciting completions reliably reflect how much a reasoning step contributes to solving the problem. The paper does not examine how these estimates change with the continuation length or number of samples. Estimating progress requires additional generations from multiple prefixes of each eligible trajectory. The paper does not report the full cost of this procedure, so training each method for the same number of epochs does not establish equal resource use.
\par\Needspace{6\baselineskip}
\begin{responsebox}[frametitle={Response to W3}]
% Edit the author response below. Reviewer text above is unchanged.
We thank the reviewer for raising these concerns. Our answer-eliciting prompt (Appendix) explicitly instructs the model to answer directly based on the current prefix, rather than continue its reasoning. Thus, the 300-token limit is intended to elicit the answer, not a full reasoning continuation. Due to time constraints, we analyzed six sampled batches. PR was evaluated for only 103/192 responses (53.6\%), as responses without successful segmentation or explicit ground-truth answer matches bypass PR estimation and remain unfiltered. Eligible responses averaged 214.8 short continuations (eight per non-final prefix, each capped at 300 new tokens), yielding 115.7 additional rollouts per original response overall, including padding. PR estimation and consistency checking incurred 118.4\% overhead relative to the inferred step-time baseline. We will report this breakdown and acknowledge the lack of compute-matched comparisons.
\end{responsebox}

\item The theory explains how aggregation can cancel some estimation errors and reduce variance under its assumptions. It does not address how using noisy progress scores to choose the segments affects these results.
\par\Needspace{6\baselineskip}
\begin{responsebox}[frametitle={Response to W4}]
% Edit the author response below. Reviewer text above is unchanged.
Our method does not assume noise-free progress estimates. Instead, sign-consistent merging and within-response ranking are designed to mitigate estimation noise. Empirically, Figure 5 shows that merging substantially improves PR--KL ranking agreement, supporting the effectiveness of this design.
\end{responsebox}

\item The results are reported without run-to-run variation or uncertainty estimates. This is a concern because improvements are modest for the second model pair and are not uniform across benchmarks. Repeated runs are needed to distinguish a reliable gain from variation in training and evaluation.
\par\Needspace{6\baselineskip}
\begin{responsebox}[frametitle={Response to W5},nobreak=true]
% Edit the author response below. Reviewer text above is unchanged.
\textit{[To be completed.]}\par\vspace{1em}
\end{responsebox}

\end{enumerate}

[1] Akhondzadeh et al. Reward-Gated On-Policy Distillation. arXiv:2607.04037v1, 2026; Section 3.

[2] Setlur et al. Rewarding Progress: Scaling Automated Process Verifiers for LLM Reasoning. ICLR 2025; arXiv:2410.08146, Section 3.1.

[3] Agentica. DeepScaleR-Preview-Dataset.

\subsection*{Justification of Recommendation}
See the weaknesses.

\subsection*{Specific Points of Feedback for Rebuttal}
\begin{enumerate}
\item What benefit is specific to the proposed selection rule, compared with simpler ways of using process or outcome rewards?
\par\Needspace{6\baselineskip}
\begin{responsebox}[frametitle={Response to Q1},nobreak=true]
% Edit the author response below. Reviewer text above is unchanged.
\textit{[To be completed.]}\par\vspace{1em}
\end{responsebox}

\item Do the improvements hold for larger students and non-Qwen backbones, such as Llama?Does the advantage remain when all methods are allowed longer responses?
\par\Needspace{6\baselineskip}
\begin{responsebox}[frametitle={Response to Q2},nobreak=true]
% Edit the author response below. Reviewer text above is unchanged.
\textit{[To be completed.]}\par\vspace{1em}
\end{responsebox}

\item How does R2-OPD compare with RG-OPD and the other relevant baselines on both model pairs?
\par\Needspace{6\baselineskip}
\begin{responsebox}[frametitle={Response to Q3},nobreak=true]
% Edit the author response below. Reviewer text above is unchanged.
\textit{[To be completed.]}\par\vspace{1em}
\end{responsebox}

\item How reliable are the progress estimates when the continuation length or number of samples is changed?
\par\Needspace{6\baselineskip}
\begin{responsebox}[frametitle={Response to Q4},nobreak=true]
% Edit the author response below. Reviewer text above is unchanged.
\textit{[To be completed.]}\par\vspace{1em}
\end{responsebox}

\item How does selecting segments from the same noisy scores affect the theoretical argument?
\par\Needspace{6\baselineskip}
\begin{responsebox}[frametitle={Response to Q5},nobreak=true]
% Edit the author response below. Reviewer text above is unchanged.
\textit{[To be completed.]}\par\vspace{1em}
\end{responsebox}

\end{enumerate}

\subsection*{Overall Response}
\par\Needspace{6\baselineskip}
\begin{responsebox}[frametitle={Overall Response to Reviewer 4 (hTMj)},nobreak=true]
% Edit the author response below. Reviewer text above is unchanged.
\textit{[To be completed.]}\par\vspace{1em}
\end{responsebox}



\clearpage

% Inlined snapshot: reviewers/reviewer5
% Reviewer 5; OpenReview ID: YgpD
% Original review text, followed by the author's supplied draft responses.
\section{Reviewer 5 (\texttt{YgpD})}
\label{reviewer:5}
\noindent\textbf{Rating:} 6: Marginally above acceptance threshold\par
\noindent\textbf{Confidence:} 4: The reviewer is confident but not absolutely certain that the evaluation is correct

\subsection*{Paper Summary}
This paper is about policy distillation, where a student model is trained by sampling its own responses and then pulling its token-level distribution toward a teacher model's distribution at every decoding step, using reverse KL divergence. The authors point out that this setup quietly assumes that being close to the teacher is the same thing as making good reasoning progress, and that assumption can break. A student span can depart from the teacher and still be the move that actually gets the problem solved, while a span that sounds just like the teacher can still be a wrong turn. Since standard policy distillation applies the same imitation pressure everywhere, it can end up punishing exactly the reasoning steps that are working and letting the unhelpful ones slide. To fix this, they propose R2 OPD. For each student response, they split it into reasoning segments using self-reflective discourse markers, estimate a solve probability at each segment boundary by sampling continuations and checking the final answer, and define a process reward as the change in solve probability across a segment. Adjacent segments with the same sign of process reward get merged into coarser units to reduce noise, and for each merged unit they also average the token-level KL loss. They then rank segments by process reward and by average KL loss within a response, and whenever a segment with higher progress also carries a higher KL loss than a segment with lower progress, that is treated as a local conflict. The most conflicting segments, up to some fixed masking ratio, get their distillation loss zeroed out entirely, so training simply stops pulling the student toward the teacher there instead of adding any new reward.

\subsection*{Strengths}
The motivating failure case is concrete rather than hypothetical. Figure 2 shows an actual example where a regressive reasoning step gets a smaller KL loss than the steps that correctly derive or recover the answer, so the reader can see exactly the kind of mismatch the paper is trying to fix rather than just taking the claim on faith.

The intervention itself is conservative in a good way. It only removes supervision it has identified as unreliable; it never adds a brand new reward term or changes the direction of the underlying objective. That makes it easy to reason about what could go wrong, since the worst case is that some tokens get less training signal, not that the model gets actively pushed in some new wrong direction.

There is real theoretical backing behind two of the design choices rather than just an intuition. Lemma 1 shows that summing process rewards over a merged run telescopes down to just the two endpoint estimates, so internal boundary noise cancels out. Proposition 2 gives a variance bound showing segment averaging of the KL loss shrinks noise as the segment gets longer under a fairly mild dependence assumption. That is a nicer level of rigor than most empirical filtering papers bother with.

The two central design choices, the masking ratio and the merging step, both get a real ablation, and the merging ablation comes with its own diagnostic beyond just downstream accuracy. Figure 5 tracks the actual rank agreement between process reward and KL loss during training and shows merging raises it well above the random baseline for almost the whole run, which is a direct check that the mechanism is doing what it claims rather than only checking whether the final score went up.

\subsection*{Weaknesses}
\begin{enumerate}
\item[W1] Three of the four cited recent baselines score below the untrained student itself, and this is never addressed. In Table 1, the plain Student baseline gets 29.67 avg@4 and 46.10 pass@4, but E OPD [1] gets 26.91 and 41.52, TIP OPD [2] gets 25.80 and 38.72, and IW OPD [3] gets 27.31 and 40.81, all clearly worse than doing no distillation at all. These three methods were each proposed in their own papers specifically to improve on plain on policy distillation, so it is quite surprising that in this paper's reimplementation they end up hurting the model relative to not training it. This could be a real property of these methods under this particular teacher and data setup, but it could just as easily be a sign that these baselines were not tuned or reproduced fairly, and the paper gives no discussion either way. Since R2 OPD's headline comparison rests on beating these same baselines, this is a fairly central concern rather than a side note.
\par\Needspace{6\baselineskip}
\begin{responsebox}[frametitle={Response to W1}]
% Edit the author response below. Reviewer text above is unchanged.
We hypothesize that the performance gap partly reflects differences in reasoning regimes. E-OPD and TIP-OPD preserve or prioritize uncertain tokens, which may encourage additional exploration without improving answer quality in thinking-enabled models. IW-OPD downweights later tokens based on accumulated teacher--student discrepancy, potentially suppressing useful verification and correction in long reasoning traces. These effects may be amplified by our response-length limit.
\end{responsebox}

\item[W2] The benchmarks driving the main comparisons are small, and there is no measure of statistical spread anywhere. AIME 2024 and AIME 2025 each contain only 30 problems, so avg@4 on either one is an average over a small number of underlying items, and a swing of a few accuracy points can come down to one or two problems flipping. No seeds, no repeated runs, and no confidence intervals are reported for any of the main table or ablation numbers. This matters more here than usual because some of the internal ablation swings are enormous for a 30-problem benchmark; for instance, Table 3 shows removing sign-consistent merging drops AIME 2024 accuracy from 32.50 to 17.5, a fifteen-point swing from one design change. A swing that large on a benchmark that small is at least as consistent with high variance as with a real underlying effect, and without any variance reporting there is no way to tell the two apart.
\par\Needspace{6\baselineskip}
\begin{responsebox}[frametitle={Response to W2},nobreak=true]
% Edit the author response below. Reviewer text above is unchanged.
\textit{[To be completed.]}\par\vspace{1em}
\end{responsebox}

\item[W3] The added cost of process reward estimation is never quantified. Getting a process reward for a single training response requires sampling $N_{\mathrm{eval}}$ on policy continuations from every accepted segment boundary in that response, on top of the response generation itself, which is a nontrivial multiple of the generation cost that plain on policy distillation does not pay at all, and that lighter baselines like TIP OPD or IW OPD likely do not pay either since they mostly reweight signals that are already available from the forward pass. The paper does not report wall clock time, extra rollouts per response, or any other measure of this overhead anywhere, so a reader cannot judge whether the accuracy gain is worth what looks like a meaningfully more expensive training recipe.
\par\Needspace{6\baselineskip}
\begin{responsebox}[frametitle={Response to W3}]
% Edit the author response below. Reviewer text above is unchanged.
Regarding overhead, across six sampled batches, PR was evaluated for 103/192 responses; others bypassed estimation. Eligible responses averaged 214.8 short continuations (eight per non-final prefix, capped at 300 tokens), with high KV-cache hit rates. These procedures added an estimated 118.4\% runtime over the inferred baseline. We acknowledge the lack of compute-matched comparisons and will clarify this limitation.
\end{responsebox}

\end{enumerate}

[1] Entropy-Aware On-Policy Distillation of Language Models.

[2] Tip: Token importance in on-policy distillation.

[3] On the Position Bias of On Policy Distillation.

\subsection*{Justification of Recommendation}
The mechanism itself is thoughtfully built, the theory around merging and averaging is a genuine plus, and the diagnostic in Figure 5 is good evidence that the method is doing what it says it does at the mechanism level.

What pulls me toward the reject side of borderline is that the actual headline comparison, R2 OPD beating recent baselines, rests on baseline numbers that look questionable on their face, since three of the four comparison methods score below doing no distillation at all, and that is never explained. On top of that, the benchmarks used for essentially every number in the paper are small enough, and the internal ablation swings large enough, that without any variance reporting I cannot tell how much of the story is signal versus benchmark noise. The missing cost accounting is a smaller concern but adds to the sense that the practical picture is incomplete.

\subsection*{Specific Points of Feedback for Rebuttal}
Can you say more about how E OPD, TIP OPD, and IW OPD were configured and tuned in your setup, and offer any explanation for why all three score below the untrained student baseline in Table 1, given that each was originally proposed to improve over plain on policy distillation.

\par\Needspace{6\baselineskip}
\begin{responsebox}[frametitle={Response to Q1},nobreak=true]
% Edit the author response below. Reviewer text above is unchanged.
\textit{[To be completed.]}\par\vspace{1em}
\end{responsebox}

\subsection*{Overall Response}
\par\Needspace{6\baselineskip}
\begin{responsebox}[frametitle={Overall Response to Reviewer 5 (YgpD)},nobreak=true]
% Edit the author response below. Reviewer text above is unchanged.
\textit{[To be completed.]}\par\vspace{1em}
\end{responsebox}




% Keep the AI review separate from the five official human reviewers.
\clearpage
\appendix

% Inlined snapshot: reviewers/ai-review
% AI review is kept separate from the five official human reviews.
\section{AI Review}
\label{reviewer:ai}
\noindent\textit{Separate AI-generated review from the OpenReview export.}

\subsection*{Synopsis of the paper}
The paper studies a mismatch between teacher imitation and productive reasoning in on-policy distillation. It proposes R\textsuperscript{2}-OPD, which estimates progress from student continuations, merges spans with consistent progress, and masks teacher supervision where progress and teacher--student divergence rankings conflict. The process signal acts only as a reliability filter rather than a training reward. Experiments cover mathematical reasoning benchmarks, multiple model pairs, and several OPD baselines.

\subsection*{Summary of Review}
The paper identifies a meaningful limitation of teacher-similarity supervision and proposes a distinctive span-level filtering mechanism. However, Equation (3), which supplies the central conflict score, is not a KL divergence and can be negative, undermining its interpretation as teacher--student disagreement. The evidence for sign-consistent merging does not establish more reliable progress estimation and largely shows recovery from a harmful unmerged variant. Reported gains also cannot be attributed specifically to progress-conflict filtering because supervision- and masking-matched controls are absent. Published precedents narrow the novelty to combining relative-progress estimation with span-level OPD masking.

\subsection*{Strengths}
\begin{itemize}
\item The distinction between imitation compatibility and reasoning progress is concretely motivated. Figure 2 illustrates how a productive alternative derivation can receive a larger divergence than a regressive but teacher-compatible span, making the targeted failure mode clear.
\item The process signal is cleanly separated from the optimized objective. Equations (11)--(13) use within-response ordering to gate the original distillation loss, avoiding direct addition of signals with incomparable scales and preserving a modular OPD update.
\item The empirical design has useful breadth within verifiable mathematics. Tables 1--2 cover two student--teacher pairs and several OPD variants, while Table 3 and Figures 4--5 expose substantial dependence on merging, masking strength, dataset, and metric rather than reporting only an aggregate result.
\end{itemize}

\subsection*{Weaknesses}
\begin{itemize}
\item Equation (3) does not define a KL divergence, although its segment average $\kappa_j$ drives every conflict decision. If $\alpha_t$ and $\beta_t$ are the student and teacher masses on $S_t$, the stated loss equals $D_{\mathrm{KL}}(q_t^S\parallel q_t^T) + \log(\alpha_t/\beta_t)$ and can be negative; its ordering can therefore reflect retained support mass rather than distributional disagreement, directly weakening Equations (11)--(12).
\par\Needspace{6\baselineskip}
\begin{responsebox}[frametitle={Response to W1},nobreak=true]
% Edit the author response below. Reviewer text above is unchanged.
\textit{[To be completed.]}\par\vspace{1em}
\end{responsebox}

\item The merging evidence does not establish that merging produces more accurate progress units. Figure 5 treats increased PR--KL agreement as reliability even though Equation (11) defines meaningful PR--KL disagreement as the target phenomenon; moreover, comparing Tables 1 and 3 shows that over 70\% of the reported 15.00 and 14.16 AIME gains over no-merging merely recover that variant's degradation below standard OPD. The evidence therefore shows that unmerged filtering is harmful, not that merging accurately denoises progress.
\par\Needspace{6\baselineskip}
\begin{responsebox}[frametitle={Response to W2},nobreak=true]
% Edit the author response below. Reviewer text above is unchanged.
\textit{[To be completed.]}\par\vspace{1em}
\end{responsebox}

\item The first process reward is not the claimed marginal progress because Equation (5) fixes $\hat{S}_0 = 0$ rather than estimating the prompt's solve probability. Although merged segment $j = 1$ is later excluded from conflict detection, its artificial nonnegative reward participates in sign-consistent merging and can absorb subsequent positive spans, changing boundaries, candidate counts, and masking decisions.
\par\Needspace{6\baselineskip}
\begin{responsebox}[frametitle={Response to W3},nobreak=true]
% Edit the author response below. Reviewer text above is unchanged.
\textit{[To be completed.]}\par\vspace{1em}
\end{responsebox}

\item The contribution is narrower than presented because its two principal ingredients have close published precedents. GenPRM considers the difference between rollout-derived success estimates at consecutive reasoning states (Zhao et al., 2026), while VOLD applies reward-guided masking to KL-based on-policy distillation (Bousselham et al., 2026); neither appears in the target paper's references, related work, or tables. The distinct contribution is therefore their span-level rank-conflict combination rather than either ingredient independently.
\par\Needspace{6\baselineskip}
\begin{responsebox}[frametitle={Response to W4},nobreak=true]
% Edit the author response below. Reviewer text above is unchanged.
\textit{[To be completed.]}\par\vspace{1em}
\end{responsebox}

\item The evaluation does not isolate the source of the gains. R\textsuperscript{2}-OPD uses ground-truth verifier supervision and eight additional continuations per evaluated boundary, but Tables 1--3 lack supervision- or compute-matched OPD, random masking with the same token budget, outcome-only masking, or direct process-reward filtering; additionally, removing merging changes segment counts and effective token masking.
\par\Needspace{6\baselineskip}
\begin{responsebox}[frametitle={Response to W5},nobreak=true]
% Edit the author response below. Reviewer text above is unchanged.
\textit{[To be completed.]}\par\vspace{1em}
\end{responsebox}

\item The reported margins lack uncertainty quantification. Table 1 improves over OPD by 2.51 aggregate avg@4 points and 4.46 pass@4 points, while Table 2's aggregate avg@4 gain is only 0.43; no training seeds, confidence intervals, or paired tests are reported. These margins cannot be separated from training and sampling variability as presented.
\par\Needspace{6\baselineskip}
\begin{responsebox}[frametitle={Response to W6},nobreak=true]
% Edit the author response below. Reviewer text above is unchanged.
\textit{[To be completed.]}\par\vspace{1em}
\end{responsebox}

\end{itemize}

\subsection*{Suggestions for Improvement}
\begin{itemize}
\item A compact algorithm block should precede the derivation and explicitly show trajectory eligibility, original-to-merged indexing, ranked-pair construction, inconsistency accumulation, masking-budget calculation, and token masking. This would remove the need to reconstruct the operational procedure across Equations (5)--(13).
\par\Needspace{6\baselineskip}
\begin{responsebox}[frametitle={Response to S1},nobreak=true]
% Edit the author response below. Reviewer text above is unchanged.
\textit{[To be completed.]}\par\vspace{1em}
\end{responsebox}

\item The mechanism analysis could distinguish whether gains arise from masking the higher-progress/high-KL endpoint or the lower-progress/low-KL endpoint of a violation. Equation (12) assigns the same pairwise contribution to both, although these cases imply different accounts of why supervision is unreliable.
\par\Needspace{6\baselineskip}
\begin{responsebox}[frametitle={Response to S2},nobreak=true]
% Edit the author response below. Reviewer text above is unchanged.
\textit{[To be completed.]}\par\vspace{1em}
\end{responsebox}

\item Results could be stratified by segmentation success, answer-string eligibility, truncation status, number of merged spans, and actual masked-token fraction. These quantities determine where R\textsuperscript{2}-OPD operates and would connect the motivating trajectory categories to downstream improvements.
\par\Needspace{6\baselineskip}
\begin{responsebox}[frametitle={Response to S3},nobreak=true]
% Edit the author response below. Reviewer text above is unchanged.
\textit{[To be completed.]}\par\vspace{1em}
\end{responsebox}

\item Claims should be scoped to verifiable mathematical reasoning and the tested small models. The construction requires ground-truth answers, executable verifiers, and repeated student continuations, so conclusions about reasoning generally extend beyond the demonstrated setting.
\par\Needspace{6\baselineskip}
\begin{responsebox}[frametitle={Response to S4},nobreak=true]
% Edit the author response below. Reviewer text above is unchanged.
\textit{[To be completed.]}\par\vspace{1em}
\end{responsebox}

\end{itemize}

\subsection*{Additional Details and Minor Issues}
\begin{itemize}
\item With $q = 30$, ceiling and exclusion of the first segment can mask half or more of the eligible segments when $n_i$ is small; Figure 4 should report actual masked-token fractions.
\par\Needspace{6\baselineskip}
\begin{responsebox}[frametitle={Response to M1},nobreak=true]
% Edit the author response below. Reviewer text above is unchanged.
\textit{[To be completed.]}\par\vspace{1em}
\end{responsebox}

\item Equation (5) evaluates intermediate states using 300-token answer-eliciting rollouts but the terminal state using the actual response's correctness, so $PR_M$ mixes different horizons and generation protocols.
\par\Needspace{6\baselineskip}
\begin{responsebox}[frametitle={Response to M2},nobreak=true]
% Edit the author response below. Reviewer text above is unchanged.
\textit{[To be completed.]}\par\vspace{1em}
\end{responsebox}

\item Table 3's displayed Olympiad values give 46.86\ensuremath{-}47.80=\ensuremath{-}0.94, rather than the reported \ensuremath{-}0.96, unless unrounded values were used.
\par\Needspace{6\baselineskip}
\begin{responsebox}[frametitle={Response to M3},nobreak=true]
% Edit the author response below. Reviewer text above is unchanged.
\textit{[To be completed.]}\par\vspace{1em}
\end{responsebox}

\item The aggregation producing values such as 43.44 AIME 2025 pass@4 in Table 1 and 18.75 avg@4 in Table 2 should be explained under the stated four-response protocol.
\par\Needspace{6\baselineskip}
\begin{responsebox}[frametitle={Response to M4},nobreak=true]
% Edit the author response below. Reviewer text above is unchanged.
\textit{[To be completed.]}\par\vspace{1em}
\end{responsebox}

\item Figure 1's answer-string category does not by itself establish a correct reasoning path or truncation caused by overthinking.
\par\Needspace{6\baselineskip}
\begin{responsebox}[frametitle={Response to M5},nobreak=true]
% Edit the author response below. Reviewer text above is unchanged.
\textit{[To be completed.]}\par\vspace{1em}
\end{responsebox}

\item Aggregated-process-reward notation changes around Equations (7), (11), and (12), and several local grammatical errors impede the otherwise linear presentation.
\par\Needspace{6\baselineskip}
\begin{responsebox}[frametitle={Response to M6},nobreak=true]
% Edit the author response below. Reviewer text above is unchanged.
\textit{[To be completed.]}\par\vspace{1em}
\end{responsebox}

\end{itemize}

\subsection*{References}
Bousselham, W., Kuehne, H., \& Schmid, C. (2026). VOLD: Reasoning transfer from LLMs to vision-language models via on-policy distillation. In Proceedings of the IEEE/CVF Conference on Computer Vision and Pattern Recognition (pp. 26209--26218). (openaccess.thecvf.com (\url{https://openaccess.thecvf.com/content/CVPR2026/html/Bousselham_VOLD_Reasoning_Transfer_from_LLMs_to_Vision-Language_Models_via_On-Policy_CVPR_2026_paper.html}))

Zhao, J., Liu, R., Zhang, K., Zhou, Z., Gao, J., Li, D., Lyu, J., Qian, Z., Qi, B., Li, X., \& Zhou, B. (2026). GenPRM: Scaling test-time compute of process reward models via generative reasoning. Proceedings of the AAAI Conference on Artificial Intelligence, 40(41), 34932--34940. (ojs.aaai.org (\url{https://ojs.aaai.org/index.php/AAAI/article/view/40797}))

\subsection*{Overall Response}
\par\Needspace{6\baselineskip}
\begin{responsebox}[frametitle={Overall Response to AI Review},nobreak=true]
% Edit the author response below. Reviewer text above is unchanged.
\textit{[To be completed.]}\par\vspace{1em}
\end{responsebox}



\end{document}
